\section{Two design laws and the reference-system theorem}
\label{sec:designlaws}

The classification (Section~\ref{sec:theorem}) and the two couplings
(Section~\ref{sec:couplings}) yield two design laws for a lunar reference-system standard
and a standard-level theorem that combines them. The logic is \Proven{}; the magnitudes it
rests on are the \Modelled{} couplings, labelled where they enter.

\begin{quote}
\textbf{Design law~1 (frame scale vs.\ timescale offset).} \emph{Geometric diversity
separates the lunar frame scale from the timescale offset in one-way ranging.} A
geometrically diverse multi-beacon multi-epoch one-way network makes the coupled Fisher
full-rank and the $\{s,\dtau\}$ marginal positive-definite (Section~\ref{sec:coupling1},
Table~\ref{tab:anchor}); a two-way tie or an external time tie is required \emph{only} in
the low-diversity / single-user regime, where the classical
$1\,\text{m}\leftrightarrow3.336\,\text{ns}$ near-defect is real. Two-way ranging confers no
selective datum-level advantage over an equal volume of well-distributed one-way
observations.
\end{quote}

\begin{quote}
\textbf{Design law~2 (the rate's residual gauge is definitional).} \emph{The lunar
timescale rate's residual gauge is definitional, not estimational.} The frame estimation
datum perturbs the rate by at most $\sim2\times10^{-17}$
\eqref{eq:framebound}, six orders of magnitude below the $\pm1.7\times10^{-11}$
reference-surface definitional band \eqref{eq:band}. No PNT observation reduces the
definitional freedom; a standard must \emph{fix} the reference surfaces (the potential
$W_0$, the selenoid) and the averaging window. Because the frame estimation barely perturbs
the rate, the rate is gauge-robust to how well the frame is measured and gauge-sensitive to
how the surfaces are defined.
\end{quote}

\begin{theorem}[A lunar reference \emph{system} must fix both a datum and a convention]
\label{thm:system}
A complete lunar reference-system realization must fix \emph{both}: (a)~a spatial datum---%
the six rigid-body gauge directions of Theorem~\ref{thm:classify}, pinned by external ties
as in \citep{baweja2026datum}; and (b)~a temporal convention---the offset $\dtau$, pinned by
an external time tie (Theorem~\ref{thm:lift}a), and the reference surfaces and averaging
window that fix the rate (Design law~2). Neither suffices alone:
\begin{enumerate}
\item[(i)] geometry alone (one-way ranging) cannot separate radial position from time
offset in the low-diversity regime (Section~\ref{sec:coupling1}, coupling~1);
\item[(ii)] the relativistic model alone cannot bootstrap an absolute rate without a spatial
datum---the rate tie \eqref{eq:ratetie} lifts $\dalpha$ only \emph{given} an independent
scale tie (Theorem~\ref{thm:lift}b); and
\item[(iii)] the rate's residual gauge is definitional, not observable, so no amount of data
removes it (Design law~2).
\end{enumerate}
\end{theorem}

\begin{proof}
Each clause restates a proven result. (i) is coupling~1: in the low-diversity regime the
$\{s,\dtau\}$ marginal has a small eigenvalue, so scale/radial and offset are not separated
by one-way geometry alone. (ii) is Theorem~\ref{thm:lift}b: the rate tie constrains only the
combination $\dalpha-\tfrac{\partial\alpha}{\partial s}\delta s$, which fixes $\dalpha$ iff
$\delta s$ is independently fixed. (iii) is Design law~2 with the bound \eqref{eq:framebound}
against the band \eqref{eq:band}. A realization missing (a) has an unobservable
frame; one missing the offset tie has an unobservable $\dtau$ (Theorem~\ref{thm:classify}ii);
one missing the surface/window convention has an undefined rate. Hence both a spatial datum
and a temporal convention are necessary.
\end{proof}

\paragraph{The P1\,$\to$\,P4 propagation.}
Theorem~\ref{thm:system} makes explicit how the spatial datum of \citep{baweja2026datum}
propagates into the timescale. The spatial datum defect enters the timescale rate through
the rate-tie Jacobian $\partial\alpha/\partial s$ \eqref{eq:dads}, so a residual scale-datum
error $\delta s$ maps to a rate error $(\partial\alpha/\partial s)\,\delta s$. This
\emph{datum-channel} gain---the genuine P1$\to$P4 gauge propagation---is
$\approx2.2\times10^{-18}$ for ILRF-class ties (the scale term of \eqref{eq:propagate}). It
must not be conflated with the broader frame-realization envelope
$\lesssim\!\sim\!2\times10^{-17}$ \eqref{eq:framebound}, which additionally absorbs non-datum
velocity and radial orbit-determination errors that Section~\ref{sec:couplings} shows do not
map onto the nine-vector datum. The spatial datum thus \emph{does} propagate into the
rate---closing the P1$\to$P4 loop---yet the rate is gauge-robust to it, and \emph{a
fortiori}: the true datum-channel gain is an order of magnitude below even that envelope. The
joint gauge is not the direct sum of the spatial and temporal gauges
(Theorem~\ref{thm:classify}iii), but along the one channel that couples them the coupling is
numerically negligible against the definitional band. This is the precise, quantified sense
in which the frame and the timescale are ``coupled but weakly so'': coupled in structure
(a nonzero $\Pst$ once the rate tie is admitted), weak in magnitude (six or more orders below
the definitional gauge).
